\documentclass[12pt, a4paper, twoside, openright]{report}

% Packages
\usepackage[utf8]{inputenc}
\usepackage[T1]{fontenc}
% Language support
\usepackage[ngerman]{babel}
\usepackage{lmodern}
\usepackage{geometry}
\usepackage{setspace}
\usepackage{graphicx}
\usepackage{hyperref}
\usepackage{booktabs}
\usepackage{caption}
\usepackage{fancyhdr}
\usepackage{titlesec}
\usepackage{etoolbox}
\usepackage{longtable}

% Fix for pandoc tightlist
\providecommand{\tightlist}{%
  \setlength{\itemsep}{0pt}\setlength{\parskip}{0pt}}

% Page geometry - MSB standard
\geometry{
  left=3cm,
  right=2.5cm,
  top=2.5cm,
  bottom=2.5cm
}

% Line spacing
\onehalfspacing

% Header/Footer
\pagestyle{fancy}
\fancyhf{}
\fancyhead[LE,RO]{\thepage}
\fancyhead[RE]{\leftmark}
\fancyhead[LO]{\rightmark}
\renewcommand{\headrulewidth}{0.4pt}

% Chapter format
\titleformat{\chapter}[hang]
  {\normalfont\Large\bfseries}{\thechapter}{1em}{}
\titlespacing*{\chapter}{0pt}{-20pt}{20pt}

% Language-specific labels
\newcommand{\labelauthor}{Autor}
\newcommand{\labelprogram}{Studiengang}
\newcommand{\labelsupervisor}{Betreuer}
\newcommand{\labeldate}{Datum}
\newcommand{\labelthesis}{Abschlussarbeit}
\newcommand{\labelabstract}{Zusammenfassung}
\newcommand{\labelacknowledgements}{Danksagung}

% Title page placeholders
\newcommand{\msbtitle}{Thesis Title}
\newcommand{\msbauthor}{Author Name}
\newcommand{\msbdegree}{Master of Science}
\newcommand{\msbprogram}{Business Administration}
\newcommand{\msbsupervisor}{Prof. Dr. Supervisor}
\newcommand{\msbdate}{\today}
\newcommand{\msblogo}{}

% Title page
\newcommand{\msbtitlepage}{
  \begin{titlepage}
    \centering
    \vspace*{1cm}
    
        \IfFileExists{\_extensions/msb/thesis/assets/logo.png}{
      \includegraphics[height=2cm]{\_extensions/msb/thesis/assets/logo.png}
      \vspace{1cm}
    }{}
        
    {\Large\bfseries Thesis Title\par}
    \vspace{1cm}
    {\large \labelthesis\par}
    \vspace{2cm}
    
    {\large
    \begin{tabular}{ll}
      \textbf{\labelauthor:} & Author Name \\
      \textbf{\labelprogram:} & Business Administration \\
      \textbf{\labelsupervisor:} & Prof.~Dr.~Betreuer Name \\
      \textbf{\labeldate:} & 2026-02-17
    \end{tabular}
    \par}
    
    \vfill
    {\large FH Münster\par}
    {\large Münster School of Business\par}
    
  \end{titlepage}
}

\makeatletter
\@ifpackageloaded{bookmark}{}{\usepackage{bookmark}}
\makeatother
\makeatletter
\@ifpackageloaded{caption}{}{\usepackage{caption}}
\AtBeginDocument{%
\ifdefined\contentsname
  \renewcommand*\contentsname{Table of contents}
\else
  \newcommand\contentsname{Table of contents}
\fi
\ifdefined\listfigurename
  \renewcommand*\listfigurename{List of Figures}
\else
  \newcommand\listfigurename{List of Figures}
\fi
\ifdefined\listtablename
  \renewcommand*\listtablename{List of Tables}
\else
  \newcommand\listtablename{List of Tables}
\fi
\ifdefined\figurename
  \renewcommand*\figurename{Figure}
\else
  \newcommand\figurename{Figure}
\fi
\ifdefined\tablename
  \renewcommand*\tablename{Table}
\else
  \newcommand\tablename{Table}
\fi
}
\@ifpackageloaded{float}{}{\usepackage{float}}
\floatstyle{ruled}
\@ifundefined{c@chapter}{\newfloat{codelisting}{h}{lop}}{\newfloat{codelisting}{h}{lop}[chapter]}
\floatname{codelisting}{Listing}
\newcommand*\listoflistings{\listof{codelisting}{List of Listings}}
\makeatother
\makeatletter
\makeatother
\makeatletter
\@ifpackageloaded{caption}{}{\usepackage{caption}}
\@ifpackageloaded{subcaption}{}{\usepackage{subcaption}}
\makeatother

\begin{document}

\msbtitlepage

\chapter*{\labelabstract}
Diese Arbeit untersucht\ldots{} Die Ergebnisse zeigen\ldots{} Die
Implikationen sind\ldots{}

\chapter*{\labelacknowledgements}
Ich möchte mich bei\ldots{} bedanken.

\tableofcontents
\listoffigures
\listoftables

\bookmarksetup{startatroot}

\chapter{Thesis Titel}\label{thesis-titel}

Diese Arbeit untersucht\ldots{} Die Ergebnisse zeigen\ldots{} Die
Implikationen sind\ldots{}

\hfill\break

\bookmarksetup{startatroot}

\chapter{Vorwort}\label{vorwort}

Willkommen zu Ihrer MSB Thesis.

\bookmarksetup{startatroot}

\chapter{Einleitung}\label{einleitung}

\bookmarksetup{startatroot}

\chapter{Einleitung}\label{einleitung-1}

\section{Hintergrund und Motivation}\label{hintergrund-und-motivation}

Hier kommt der Hintergrund Ihrer Arbeit\ldots{}

\section{Forschungsfrage}\label{forschungsfrage}

Die zentrale Forschungsfrage lautet: \ldots{}

\section{Aufbau der Arbeit}\label{aufbau-der-arbeit}

Die Arbeit ist wie folgt aufgebaut:

\begin{itemize}
\tightlist
\item
  Kapitel 2: Literaturübersicht
\item
  Kapitel 3: Methodik
\item
  Kapitel 4: Ergebnisse
\item
  Kapitel 5: Diskussion
\item
  Kapitel 6: Fazit
\end{itemize}

\bookmarksetup{startatroot}

\chapter{Literaturübersicht}\label{literaturuxfcbersicht}

\bookmarksetup{startatroot}

\chapter{Literaturübersicht}\label{literaturuxfcbersicht-1}

\section{Theoretischer Rahmen}\label{theoretischer-rahmen}

Dieser Abschnitt stellt den theoretischen Rahmen dar\ldots{}

\section{Stand der Forschung}\label{stand-der-forschung}

Der aktuelle Stand der Forschung zeigt\ldots{} Smith (2020)
argumentiert, dass\ldots{}

\section{Forschungslücke}\label{forschungsluxfccke}

Es besteht eine Forschungslücke in\ldots{}

\bookmarksetup{startatroot}

\chapter{Methodik}\label{methodik}

\bookmarksetup{startatroot}

\chapter{Methodik}\label{methodik-1}

\section{Forschungsdesign}\label{forschungsdesign}

Diese Studie verwendet ein quantitatives Forschungsdesign\ldots{}

\section{Datenerhebung}\label{datenerhebung}

Die Daten wurden erhoben durch\ldots{}

\section{Datenanalyse}\label{datenanalyse}

Zur Analyse wurden folgende Methoden angewendet:

\begin{itemize}
\tightlist
\item
  Deskriptive Statistik
\item
  Regressionsanalyse
\item
  Hypothesentests
\end{itemize}

\bookmarksetup{startatroot}

\chapter{Ergebnisse}\label{ergebnisse}

\bookmarksetup{startatroot}

\chapter{Ergebnisse}\label{ergebnisse-1}

\section{Deskriptive Statistik}\label{deskriptive-statistik}

Die deskriptiven Ergebnisse zeigen\ldots{}

Siehe Abbildung 1 für eine Visualisierung der Hauptergebnisse.

\section{Hypothesentests}\label{hypothesentests}

Die Ergebnisse der Hypothesentests sind in Table~\ref{tbl-results}
dargestellt.

\begin{longtable}[]{@{}lll@{}}
\caption{Hypothesentestergebnisse}\label{tbl-results}\tabularnewline
\toprule\noalign{}
Hypothese & p-Wert & Ergebnis \\
\midrule\noalign{}
\endfirsthead
\toprule\noalign{}
Hypothese & p-Wert & Ergebnis \\
\midrule\noalign{}
\endhead
\bottomrule\noalign{}
\endlastfoot
H1 & 0.001 & Bestätigt \\
H2 & 0.123 & Abgelehnt \\
\end{longtable}

\bookmarksetup{startatroot}

\chapter{Diskussion}\label{diskussion}

\bookmarksetup{startatroot}

\chapter{Diskussion}\label{diskussion-1}

\section{Interpretation der
Ergebnisse}\label{interpretation-der-ergebnisse}

Die Ergebnisse zeigen, dass\ldots{}

\section{Theoretische Implikationen}\label{theoretische-implikationen}

Aus theoretischer Sicht\ldots{}

\section{Praktische Implikationen}\label{praktische-implikationen}

Für die Praxis bedeutet dies\ldots{}

\section{Limitationen}\label{limitationen}

Diese Studie unterliegt folgenden Limitationen\ldots{}

\bookmarksetup{startatroot}

\chapter{Fazit}\label{fazit}

\bookmarksetup{startatroot}

\chapter{Fazit}\label{fazit-1}

\section{Zusammenfassung}\label{zusammenfassung}

Diese Arbeit hat gezeigt, dass\ldots{}

\section{Ausblick}\label{ausblick}

Zukünftige Forschung sollte\ldots{}


\end{document}
